\documentclass[11pt]{article}

\usepackage[a4paper, margin=27mm]{geometry}
\usepackage{mathptmx}            % Times — the Bitcoin-paper look
\usepackage[T1]{fontenc}
\usepackage[utf8]{inputenc}
\usepackage{microtype}
\usepackage{booktabs}
\usepackage{amsmath}
\usepackage{enumitem}
\usepackage{listings}
\usepackage{xcolor}
\usepackage{textcomp}
\usepackage[hidelinks]{hyperref}
\usepackage{parskip}

\setlength{\parskip}{4pt plus 1pt}
\setlength{\parindent}{0pt}

\lstset{
  basicstyle=\footnotesize\ttfamily,
  frame=single,
  framerule=0.3pt,
  rulecolor=\color{black!40},
  backgroundcolor=\color{black!4},
  xleftmargin=8pt, xrightmargin=8pt,
  aboveskip=8pt, belowskip=8pt,
  breaklines=true,
  columns=fullflexible,
  keepspaces=true,
}

% Listings wrapped in a minipage cannot break across pages
\lstnewenvironment{code}
  {\medskip\noindent\minipage{\linewidth}}
  {\endminipage\medskip}

\begin{document}

\begin{center}
{\Large\bfseries Freeport: A Peer-to-Peer Marketplace Protocol}\\[10pt]
{\normalsize Phil T}\\[2pt]
{\small\texttt{ptrinh@me.com}}\\[8pt]
{\small Draft v0.2 --- June 2026}
\end{center}

\vspace{14pt}

\begin{center}
\begin{minipage}{0.88\textwidth}
\small
\textbf{Abstract.} A purely peer-to-peer marketplace would allow buyers and
sellers of local goods and services to transact directly without going
through a platform intermediary. Digital signatures and open relays provide
part of the solution, but the main benefits are lost if a trusted third
party is still required to establish reputation. We propose a solution to
the trust problem using a web of co-signed deal receipts. Participants
publish signed \emph{intents} to open relays; counterparties negotiate terms
over encrypted messages using a bounded state machine; a deal is proven only
when both parties independently sign a receipt referencing it. Reputation is
not a global score but a subjective quantity each viewer computes by walking
the receipt graph outward from their own key, so fabricated identities ---
however numerous --- carry no weight with anyone they have not actually
served. The protocol requires no platform, no fees, no account, and no
permission. Markets are defined by convention: any group of people agreeing
on a topic tag and a payload schema constitutes a marketplace.
\end{minipage}
\end{center}

\section{Introduction}

Commerce between individuals on the internet has come to rely almost
exclusively on platforms serving as trusted third parties: ride-hailing
companies, gig marketplaces, classified-ad sites. While the model works well
enough for most transactions, it suffers from the inherent weaknesses of the
trust-based model. Platforms extract commissions of 15--30\% of transaction
value. They unilaterally set terms, suspend accounts without recourse, and
hold reputation hostage --- a seller's years of reviews cannot leave the
platform that hosts them. Network effects entrench incumbents; participants
on both sides pay monopoly rents. The platform mediates disputes it has
incentives to resolve in its own interest, and it surveils every transaction
as a condition of participation.

What is needed is a marketplace based on cryptographic proof instead of
platform trust, allowing any two willing parties to find each other, agree
on terms, and build durable reputations without a trusted third party. In
this paper, we propose a solution using signed intents on open relays,
encrypted bilateral negotiation, and a reputation system in which proof of
past trade is established by mutual digital signature and weighted by each
observer's own position in the trade graph.

\section{Identities}

An identity is a secp256k1 key pair, generated locally on first use, as in
Nostr~[2]. There is no registration, no account, and no issuing authority.
The public key identifies a participant across all markets; the private key
signs every event the participant publishes and decrypts messages addressed
to them. Identity is portable by construction: reputation accrues to the
key, not to any service, and the key can be backed up (encrypted under a
passphrase per NIP-49) and restored on any client.

Keys are free, which makes identity creation costless --- the central
problem this design must address. We return to it in
Sections~\ref{sec:subjective}--\ref{sec:privacy}.

\section{Intents}

A market posting is an \emph{intent}: a signed, addressable event declaring
that its author wants to buy or sell something, where, when, and on what
terms.

\begin{code}
kind:  32101 (offer) | 32102 (request)
tags:  d = stable intent id          (replaceable by author)
       t = market topic              (e.g. "sg-rideshare")
       g = geohash prefix            (location scope)
       expiration                    (NIP-40)
content: { schema, title, payload, window, flex_minutes, ... }
\end{code}

The \texttt{payload} is schema-keyed (\texttt{rideshare/1},
\texttt{service/1}, \dots) so the protocol is vertical-agnostic: a rideshare
market, a plumbing market, and a homemade-food market differ only in payload
schema and topic tag. Anyone can define a new market by publishing intents
under a new tag; no registry exists and none is needed. Geohash tags scope
discovery to a neighborhood; expiration tags let relays garbage-collect
stale postings. Because intents are addressable, an author can revise or
withdraw a posting by republishing under the same \texttt{d} tag --- a
withdrawal is just a republished event carrying a near-term expiration, so
relays drop it on schedule rather than serving a posting the author has
abandoned.

Each intent carries a small proof-of-work (NIP-13) so that posting into a
market has a nonzero cost, throttling flooding without a gatekeeper. The work
is modest --- a few leading zero bits --- and clients mine it off the main
thread so the interface stays responsive while it is computed.

Relays are the only infrastructure, and they are interchangeable commodity
infrastructure: dumb stores that accept signed events and serve filtered
subscriptions~[2]. A market lives on whichever relays its participants
choose to use; censorship by one relay is routed around by publishing to
others. Clients publish to several relays and read from several, so no relay
is a point of control.

\section{Negotiation}

Discovery is public; negotiation is private. When a counterparty responds to
an intent, the parties exchange encrypted direct messages (NIP-04, with
NIP-17 planned) carrying JSON envelopes that drive a small state machine:

\begin{code}
open --counter--> open ... (<= 8 rounds)
  |
accept --> accepted_by_us / accepted_by_them
  |                          |
  +----- second accept ------+--> confirmed
  |                                  |
  |               owner cancel ------+ (filled) --> cancelled
cancel / timeout --> cancelled / expired
\end{code}

Each \texttt{counter} proposes a complete replacement set of terms (time
window, payment, route or scope), so there is no ambiguity about what is on
the table. A deal requires a \emph{double accept}: one party accepts the
standing terms, and the deal is confirmed only when the other party accepts
in turn. This closes the race in which both sides simultaneously counter and
each believes a different version was agreed. The round limit bounds
negotiation spam.

A frequent special case collapses the double accept to a single tap. When a
responder accepts the standing terms \emph{unchanged} --- a driver taking a
ride at the posted price --- no new terms are introduced, so the client
records an immediate confirmation rather than opening another round.
Confirmation remains provisional on the responder's side, however, because a
single request can draw several competing responders. The intent owner
confirms at most one: on doing so the owner's client withdraws the intent and
sends an authoritative cancellation --- \emph{filled, taken by another offer}
--- to the losing responders. Receiving it moves a negotiation the responder
had locally marked \texttt{confirmed} to \texttt{cancelled}, resolving the
race so that no two parties believe they hold the same deal. A
\texttt{completed} deal is terminal and is never reversed by such a message.

On confirmation, the parties exchange contact information inside the encrypted
channel --- by design the protocol's job ends at the handshake; the ride or
the repair happens in the physical world.

\section{Deal Receipts}

The negotiation transcript is private and either party could fabricate one,
so confirmed deals must be anchored publicly. When a negotiation reaches
\texttt{confirmed}, each party publishes a \emph{receipt}:

\begin{code}
kind:  32104
tags:  d = negotiation id, p = counterparty
\end{code}

A deal is \textbf{proven} if and only if both halves exist: $A$ signed a
receipt naming $B$, and $B$ signed a receipt naming $A$, for the same
negotiation id. A single key cannot manufacture a proven deal; it takes two
signatures from two keys. This does not prevent one person from controlling
both keys --- no cryptographic scheme can --- but it forces every fabricated
deal to construct an edge in a public graph, and
Section~\ref{sec:subjective} shows why fabricated edges are worthless.

Receipts are addressable by negotiation id, so each party can publish at
most one receipt per deal, and they reference the intent so the full chain
--- posting, negotiation id, mutual confirmation --- is auditable by any
verifier.

\section{Ratings}

After a proven deal, each party may rate the other once:

\begin{code}
kind:  32103
tags:  d = negotiation id, p = ratee
content: { score in {-1, 0, +1, +2}, note?,
           contact_verified?, contact_masked? }
\end{code}

The d-tag makes ratings idempotent per (rater, deal) --- re-rating replaces
rather than accumulates. A verifier counts a rating only if (a)~the
negotiation id matches a proven receipt pair and (b)~the rater is the
counterparty of that very deal. Ratings detached from proven deals are
discarded. Repeated ratings between the same pair decay geometrically
($1, \tfrac{1}{2}, \tfrac{1}{4}, \dots$): the quantity that matters is
\emph{distinct counterparties}, which is exactly the quantity made expensive
by Section~\ref{sec:subjective}.

The optional \texttt{contact\_verified} field is a peer attestation that the
rater actually reached the counterparty at their advertised (masked) phone
number, binding the public contact claim to a private exchange --- see
Section~\ref{sec:privacy}. Rating events additionally carry a small
proof-of-work (NIP-13) as a spam floor.

\section{Reputation as a Subjective Quantity}
\label{sec:subjective}

It is impossible to build a Sybil-proof global reputation score when
identities are free~[4]. Rather than fight this theorem, the protocol
abandons the global score. Reputation in Freeport is \emph{subjective}: each
viewer computes the weight of every rater from their own vantage point by
breadth-first traversal of the proven-receipt graph:

\begin{center}
\begin{tabular}{lll}
\toprule
hop 0 & people the viewer has proven deals with & weight 1.0 \\
hop 1 & their proven counterparties             & weight 0.3 \\
hop 2 & one ring further                        & weight 0.1 \\
---   & unreachable                             & weight $\approx 0$ \\
\bottomrule
\end{tabular}
\end{center}

The viewer's social follow list (NIP-02) seeds the map so that newcomers
with no deal history still see through the eyes of people they chose to
follow. A subject's displayed reputation is then the trust-weighted average
of valid ratings, alongside the raw counts: \emph{deals, distinct partners,
partners within your network}.

Consider an attacker who fabricates $n$ puppet identities and has them trade
and rate each other indefinitely. Every rating passes the receipt check ---
the attacker controls both signatures. But the puppet cluster is connected
to an honest viewer only if some honest participant has a proven deal into
it. With no such edge, every puppet rater carries the unreachable-rater
floor weight $\varepsilon$, and the attacker's influence on the viewer's
perception is bounded by $\varepsilon$ regardless of $n$. The attack
surfaces in the UI as its own signature:

\begin{code}
honest:  * 12 deals - 9 partners - 3 in your network
           - verified by 9
sybil:   50 deals - 50 partners - 0 in your network
           - new account
\end{code}

To gain real influence the attacker must complete genuine deals with
genuinely connected participants --- at which point he is simply a
participant. The cost of attacking the system converges to the cost of using
it honestly, which is the property we require. Dampeners narrow the
remaining margins: a rater whose entire receipt history is deals with one
subject is discounted (the two-puppet pattern), accounts with no event
history older than a week are flagged, and unknown-rater weight
$\varepsilon$ keeps distant ratings visible but uninfluential.

\section{Privacy and Contact Binding}
\label{sec:privacy}

Identity is pseudonymous by default; the protocol adds disclosure only where
it pays for trust, and lets the user choose the dose.

\textbf{Phone numbers.} A participant may attach a phone number to their
profile, published either in full (explicit opt-in, with a permanence
warning) or masked to country code and final digits
(\texttt{+84\textbullet\textbullet\textbullet\textbullet\textbullet\textbullet 678}).
Masking is applied \emph{before} publication; the full number never leaves
the device except inside the encrypted channel at deal confirmation. A
self-published mask proves nothing --- anyone can publish any string --- so
the protocol binds it to reality through peers: at deal time the full number
travels to the counterparty over the encrypted channel; when rating, the
counterparty's client attests the canonical mask of the number it
\emph{actually reached}. Verifiers cross-check attested masks against the
subject's published mask and discard attestations on mismatch. A fake public
mask is thus exposed by the first honest deal, and ``verified by $N$
partners'' means $N$ proven counterparties independently confirmed
reachability --- a peer-produced verification requiring no SMS provider and
no central verifier.

\textbf{Negotiations} are end-to-end encrypted and invisible to relays
beyond traffic metadata. \textbf{Location} is published only at geohash
precision chosen by the client ($\approx \pm 0.6$\,km), honest about its own
imprecision. Reputation binds to the key, not the phone: rotating numbers
neither transfers nor escapes a rating history.

\section{Incentives}

The protocol has no token, no fee, and no operator, which raises the
question of why infrastructure exists. The answer is that the infrastructure
is nearly free. Relays are commodity event stores, already operated in the
hundreds by the Nostr ecosystem for social traffic; a community marketplace
fits comfortably in their economics, and any community can run its own relay
for the cost of a small server. Clients are ordinary apps competing on
quality, none privileged by the protocol. The parties to a deal settle
payment by whatever means they already trust --- cash, bank transfer, or, in
a later protocol phase, Lightning~[5], whose escrow-by-hold-invoice could
also serve as a costly-to-forge signal strengthening
Section~\ref{sec:subjective}'s graph.

What the design deliberately forecloses is the platform's rent: there is no
position in the protocol from which to take a cut of trades, because there
is no chokepoint through which trades must pass. Discovery, negotiation,
proof, and reputation are all client-side computations over public, signed
data.

\section{Comparison with Platforms}

\begin{center}
\begin{tabular}{lll}
\toprule
 & \textbf{Platform} & \textbf{Freeport} \\
\midrule
Listing    & account required, ToS      & signed event, no permission \\
Matching   & proprietary algorithm      & open filters, client-side \\
Fees       & 15--30\% commission        & none \\
Reputation & custodial, non-portable    & signed, portable, subjective \\
Dispute    & platform adjudication      & reputation consequences \\
Censorship & account suspension         & publish to another relay \\
Privacy    & full transaction surveillance & encrypted negotiation, pseudonyms \\
\bottomrule
\end{tabular}
\end{center}

The honest trade-off is dispute resolution: a platform can refund a buyer
and ban a seller; Freeport can only ensure the bad deal is provably
attributable and permanently visible to every future counterparty who looks.
For low-value, high-frequency local commerce --- rides, repairs, food ---
reputation consequences are the mechanism that already governs offline
behavior, and the protocol's contribution is to make that mechanism
portable, verifiable, and unforgeable at scale.

\section{Conclusion}

We have proposed a marketplace without a trusted third party. Participants
declare intents on open relays, negotiate privately under a bounded state
machine, and anchor completed deals with mutual signatures. Reputation is
recast from a global score --- impossible to defend with free identities ---
into a subjective quantity computed over the proven-deal graph from each
viewer's own position, making Sybil clusters visible and weightless while
honest history compounds. Contact claims are bound to reality by peer
attestation rather than by any verifier. The protocol is small: four event
kinds and one message envelope over an existing relay network. Markets
themselves are conventions --- a tag and a schema --- so anyone may open
one, and no one may close one.

\begin{thebibliography}{9}
\small

\bibitem{nakamoto} S.~Nakamoto, ``Bitcoin: A Peer-to-Peer Electronic Cash
System,'' 2008.

\bibitem{nostr} fiatjaf et al., ``Nostr: Notes and Other Stuff Transmitted
by Relays --- NIP-01: Basic protocol flow,''
\texttt{github.com/nostr-protocol/nips}.

\bibitem{nips} Nostr NIPs employed: NIP-02 (follow lists), NIP-04/NIP-17
(encrypted messaging), NIP-13 (proof of work), NIP-19 (bech32 entities),
NIP-40 (expiration), NIP-44 (versioned encryption), NIP-49 (key encryption),
NIP-96/98 (HTTP file storage and auth).

\bibitem{douceur} J.~R.~Douceur, ``The Sybil Attack,'' \emph{IPTPS}, 2002.

\bibitem{lightning} J.~Poon, T.~Dryja, ``The Bitcoin Lightning Network:
Scalable Off-Chain Instant Payments,'' 2016.

\bibitem{sybilrank} Q.~Cao, M.~Sirivianos, X.~Yang, T.~Pregueiro, ``Aiding
the Detection of Fake Accounts in Large Scale Social Online Services''
(SybilRank), \emph{NSDI}, 2012.

\bibitem{sybilguard} H.~Yu, M.~Kaminsky, P.~B.~Gibbons, A.~Flaxman,
``SybilGuard: Defending Against Sybil Attacks via Social Networks,''
\emph{SIGCOMM}, 2006.

\end{thebibliography}

\end{document}
